% ============================================================
%  Chapter 7, User Documentation
% ============================================================
\chapter{User Documentation}
\label{ch:user_docs}

\section{Use Cases}
\label{sec:usecases}

The system distinguishes two actors. The \textbf{User} (researcher, developer)
operates the trading engine and the web dashboard, while \textbf{Binance}
acts as an external system providing market data and accepting orders.
Figure~\ref{fig:usecases} summarises the supported use cases.

\begin{figure}[H]
  \centering
  \resizebox{\textwidth}{!}{%
  \begin{tikzpicture}[
      every node/.style={font=\small},
      actor/.style={draw, thick, minimum width=2.2cm, align=center},
      uc/.style={draw, ellipse, thick, minimum width=3.6cm, minimum height=0.9cm, align=center, inner sep=2pt}
    ]
    \node[actor] (user) at (0,0) {\textbf{User}\\(researcher, developer)};
    \node[actor] (binance) at (12,0) {\textbf{Binance}\\(external system)};

    \node[uc] (uc1) at (6, 3.5)  {Run backtest};
    \node[uc] (uc2) at (6, 2.2)  {Start paper trading};
    \node[uc] (uc3) at (6, 0.9)  {Start live trading};
    \node[uc] (uc4) at (6,-0.4)  {Train model offline};
    \node[uc] (uc5) at (6,-1.7)  {Inspect dashboard};
    \node[uc] (uc6) at (6,-3.0)  {Configure risk policy};

    \node[draw, thick, dashed, fit=(uc1)(uc2)(uc3)(uc4)(uc5)(uc6),
          inner sep=10pt, label=above:{\textbf{KCMamba trading system}}] {};

    \foreach \u in {uc1,uc2,uc3,uc4,uc5,uc6} {
      \draw (user) -- (\u);
    }
    \draw (uc2) -- (binance);
    \draw (uc3) -- (binance);
  \end{tikzpicture}%
  }
  \caption{Use case diagram of the KCMamba system. The user has six main
           operations available; Binance participates as an external system
           on the paper- and live-trading paths (market data, order
           execution).}
  \label{fig:usecases}
\end{figure}

\section{Purpose and Audience}

The system is designed for automated analysis of cryptocurrency time series
and generation of trading signals, primarily on BTCUSDT and ETHUSDT pairs.
Its target audience is developers and researchers who want to backtest or
live-run ML-based trading strategies.

\section{Overview of the Methods Used}

The system relies on the \textbf{KCMamba} architecture for price forecasting,
which embeds a differentiable Kalman filter layer into the Mamba selective
state space model (SSM): learned $Q_\text{net}$, $R_\text{net}$, and
$K_\text{net}$ subnetworks estimate the per-step ratio of process and
measurement noise, and the resulting adaptive $K$-gain controls how much
the model trusts the current observation versus the prior state (full
description in Chapter~\ref{ch:architektura}). The input OHLCV candles are
turned into sliding-window feature vectors, from which the model produces a
96-step forecast based on a 96-step lookback. A threshold-based strategy
(\texttt{ThresholdStrategy}) converts the forecasts into
\textsc{Buy}/\textsc{Sell}/\textsc{Hold} signals, which a risk-management
layer (position-size cap, stop-loss, daily loss limit) filters before
the order is dispatched through the Binance Futures REST endpoint. The
architectural decisions, training protocol, and experimental evaluation are
detailed in Chapters~\ref{ch:architektura} and~\ref{ch:kiserletek}.

\section{System Requirements}

\begin{table}[H]
  \centering
  \small
  \begin{tabular}{|l|l|l|}
    \hline
    \textbf{Component} & \textbf{Minimum} & \textbf{Recommended} \\
    \hline\hline
    Operating system & Linux/macOS/Windows 10+ & Ubuntu 22.04 LTS \\
    Python           & 3.10                   & 3.11 \\
    RAM              & 4~GB                   & 8~GB \\
    Disk             & 2~GB                   & 5~GB \\
    GPU (optional)   & NVIDIA $\ge$6~GB VRAM  & NVIDIA $\ge$12~GB VRAM \\
    Network          & n.a. (backtest mode)   & Stable internet (live) \\
    \hline
  \end{tabular}
  \caption{Minimum and recommended hardware/software requirements.}
  \label{tab:sysreq}
\end{table}

\section{Installation Guide}

\begin{enumerate}
  \item Clone and create a virtual environment:
\begin{verbatim}
git clone <repo_url> && cd diplomamunkakod
python -m venv .venv && source .venv/bin/activate
\end{verbatim}
  \item Install dependencies (ML packages are optional):
\begin{verbatim}
pip install -r requirements.txt
pip install -r requirements-ml.txt   # GPU: CUDA-compatible PyTorch
\end{verbatim}
  \item Configure API keys (only required for live/paper mode):
\begin{verbatim}
cp .env.example .env
# Edit .env: BINANCE_API_KEY, BINANCE_SECRET
\end{verbatim}
\end{enumerate}

\section{Usage Guide}

\subsubsection*{First run (paper-trading demo)}

The default way to run the system is the foreground paper-trading demo.
Executing \texttt{main.py} launches three background services concurrently,
namely the offline trainer, the Binance WebSocket feed, and the web dashboard.

\begin{verbatim}
source .venv/bin/activate
python main.py
\end{verbatim}

On a successful start, the console prints the following blocks in order: the
\texttt{core} module imports, adapter initialisation, a historical warm-up
download (by default 3 days of BTCUSDT 5-minute candles, controlled by
\texttt{warmup\_start\_str}), and finally the
\texttt{[BOOT] TradingEngine started} banner, after which the dashboard is
available at \texttt{http://127.0.0.1:8080}.

\subsubsection*{Dashboard views}

The top navigation bar exposes eight views grouped into three categories:

\begin{table}[H]
  \centering
  \small
  \begin{tabular}{|l|l|p{7.2cm}|}
    \hline
    \textbf{Category} & \textbf{Page} & \textbf{Content} \\
    \hline\hline
    Financial & \texttt{/}              & Equity curve, open positions, live price chart \\
    Financial & \texttt{/signals.html}  & Generated signal stream with risk status \\
    Financial & \texttt{/model.html}    & Model version, calibration and diagnostic metrics \\
    Financial & \texttt{/training.html} & Offline training progress, loss curves \\
    Benchmark & \texttt{/benchmark.html}     & Synthetic noisy benchmark results \\
    Benchmark & \texttt{/ltsf\_benchmark.html}& LTSF (Exchange, ETTh) benchmark results \\
    KCA analysis & \texttt{/kla\_kalman.html}   & Live Kalman parameter ($K$, $A$, $R$) traces \\
    KCA analysis & \texttt{/complexity.html}    & Parameter count $\times$ performance heatmap \\
    \hline
  \end{tabular}
  \caption{Role of each dashboard view.}
  \label{tab:dashboard_pages}
\end{table}

A context-sensitive help dialog opens on clicking the \texttt{?} button in the
bottom-right corner of every page. Page-to-page navigation happens through the
top bar, and the browser's \emph{back} button also works.

\subsubsection*{Usage patterns}

The system supports four typical scenarios. The paper-trading demo (default)
runs on the Binance live feed with simulated execution. The backtest mode
runs the same pipeline against historical data without a Binance connection
via \texttt{adapters/testing/backtest.py}, writing results to the console and
CSV. The offline training mode ingests a full dataset (by default 3 days of
BTCUSDT 5-minute candles) in a single pass and persists the weights to a
\texttt{*.pt} file. Finally, a \emph{dashboard-only} mode is useful for
reviewing a previous run: the Binance feed can be disabled and the dashboard
browsed offline.

\subsubsection*{Runtime logging and shutdown}

\texttt{stdout} contains normal status messages
(\texttt{[BOOT]}, \texttt{[SIGNAL]}, \texttt{[FILL]}, \texttt{[DASHBOARD]}),
while \texttt{stderr} holds warnings and errors
(\texttt{[WARN]}, \texttt{[FATAL ERROR]}). The two streams can be logged
separately:
\begin{verbatim}
nohup python main.py > run_log.txt 2> run_err.txt &
\end{verbatim}

The run can be safely interrupted at any time with \textbf{Ctrl+C}. The
\texttt{finally} clause guarantees that
(i) the WebSocket connection is closed,
(ii) in-memory event buffers are flushed,
(iii) the HTTP server closes active connections, and
(iv) the final state (open positions, equity, last signal) is written under
the \texttt{FINAL DASHBOARD SNAPSHOT} header.

\subsubsection*{Security considerations}

The demo runs in \emph{paper-trading} mode by default and does not submit
real Binance orders. If a production key is nonetheless placed in the
\texttt{.env} file, it is recommended to (a) grant only read and Futures
trading permissions on the Binance side, and (b) restrict the key by
IP whitelist. The dashboard listens on \texttt{0.0.0.0:8080} by default, and
exposing it on a public IP should only be done behind an external
reverse proxy (nginx, traefik).

\section{Troubleshooting}

\noindent\textbf{\texttt{ModuleNotFoundError: mamba\_ssm}}\\
Install the CUDA-compatible packages from \texttt{requirements-ml.txt}, or
run on CPU: \texttt{python main.py --device cpu}.

\noindent\textbf{Binance API 403/401}\\
Check the \texttt{.env} keys and IP whitelist settings on the Binance trading
interface.

\noindent\textbf{WebSocket connection error}\\
A firewall may block Binance's \texttt{wss://} endpoint (port $9443$). Set
up a proxy if necessary.

\section{Runtime Messages}

The table below summarises the most common runtime messages and
recommended actions.

\begin{table}[H]
  \centering
  \small
  \begin{tabular}{|p{4.4cm}|p{2.8cm}|p{6.2cm}|}
    \hline
    \textbf{Message / error} & \textbf{Severity} & \textbf{Action} \\
    \hline\hline
    \texttt{[WARN] Offline training failed} & Warning & System continues with the previous model; check network access and data availability. \\
    \hline
    \texttt{[WebSocket] Connecting...} & Info & Automatic reconnect; if it fails after 5 attempts, check firewall (port $9443$). \\
    \hline
    \texttt{ModuleNotFoundError: mamba\_ssm} & Error & Install \texttt{requirements-ml.txt} packages or pass \texttt{--device cpu}. \\
    \hline
    \texttt{Binance API 401/403} & Error & Invalid API key; update \texttt{.env} and verify IP whitelist on Binance. \\
    \hline
    \texttt{Not enough data (N rows)} & Warning & Extend \texttt{training\_start\_str} window or wait for more closed candles. \\
    \hline
    \texttt{[FATAL ERROR] ...} & Critical & Stack trace on \texttt{stderr}; send the full log to the developer. \\
    \hline
  \end{tabular}
  \caption{Common runtime messages and recommended troubleshooting steps.}
  \label{tab:errortable}
\end{table}
